\section{Reverse Words in a String}
\label{sec:str-reverse-words}

\problemstatement{Given a string of words separated by spaces (possibly with
leading, trailing, or multiple spaces), return the words in reverse order,
separated by single spaces.}

\begin{patternbox}
\emph{``Reverse the order of words''} is a parse-and-rebuild: extract words
by scanning, then emit them back-to-front. Handling stray spaces is the only
subtlety, cleanly solved with a string stream or a manual scan.
\end{patternbox}

\subsection*{Split, then join in reverse}

Scan the string, collecting maximal non-space runs as words while skipping
any spaces. Store the words, then output them from last to first joined by a
single space. This normalises whitespace automatically -- extra and edge
spaces vanish because only actual words are collected.

\begin{ideabox}[Collect Words, Ignore the Spaces]
Treating the string as a sequence of words rather than characters sidesteps
messy space handling: gather each word once, discard the gaps, and reorder.
The single-space join produces clean output regardless of the input's
spacing.
\end{ideabox}

\subsection*{Algorithm}

\begin{enumerate}
  \item Read words one at a time (a stream skips runs of spaces).
  \item Push each word; then emit them in reverse with single-space
        separators.
\end{enumerate}

\subsection*{Dry run}

\dryrun{\texttt{"  the sky\ \ is  blue "}.}

\begin{itemize}
  \item Words: \texttt{the}, \texttt{sky}, \texttt{is}, \texttt{blue}.
  \item Reversed and joined: \texttt{"blue is sky the"}.
\end{itemize}

\subsection*{C++ implementation}

\begin{lstlisting}
#include <bits/stdc++.h>
using namespace std;

string reverseWords(const string& s) {
    stringstream ss(s);
    string word;
    vector<string> words;
    while (ss >> word) words.push_back(word);      // skips extra spaces
    string res;
    for (int i = (int)words.size() - 1; i >= 0; i--) {
        res += words[i];
        if (i) res += ' ';
    }
    return res;
}

int main() {
    cout << "[" << reverseWords("  the sky is  blue ") << "]\n";  // [blue is sky the]
    return 0;
}
\end{lstlisting}

\begin{complexitybox}
\textbf{Time Complexity.} $O(n)$. \textbf{Space Complexity.} $O(n)$ for the
words and output.
\end{complexitybox}

\begin{takeawaybox}
Word-level problems become easy when you tokenise first and treat words as
units; a stringstream handles arbitrary whitespace for free. In-place
$O(1)$-space variants reverse the whole string then each word.
\end{takeawaybox}
